% =====================================================================
% Section 6 — Conclusion
%
% Written only after Sections 4 and 5 were complete and every table and
% figure was frozen. Introduces no number that does not already appear in
% Section 4, and no claim not already argued in Section 5.
% =====================================================================

\section{Conclusion}
\label{sec:conclusion}

This study benchmarked five deep-learning architectures on macroscopic
cross-section identification of 17 conifer species, under a partition
grouped by origin key and evaluated from a single immutable registry.

Three results are supported by the evidence presented. First, ResNet-50 is
not competitive on this task: at $92.356 \pm 1.022$\,\% accuracy it sits
roughly seven points below ConvNeXt-Tiny, a separation several times its own
seed-to-seed variation, and it accounts for 122 of the 204 errors made by
all five backbones across all seeds. Second, the remaining four architectures
are not separated by this experiment. ConvNeXt-Tiny, Swin-Tiny and
MobileNetV3-Large differ by margins comparable to the standard deviation
across three seeds, so their ordering should not be treated as established.
Third, within that group MobileNetV3-Large is the practically interesting
result: it comes within 1.19 accuracy points of the best model using one
sixth the parameters and one twentieth the compute, which is the trade-off
that matters for field deployment.

The error structure is concentrated rather than diffuse. In the strongest
model every error across three seeds is the same confusion, \textit{Pinus
krempfii} predicted as \textit{Pinus dalatensis}, and that pair is the only
one appearing in all five architectures. It is strongly asymmetric, 61
occurrences against one in the reverse direction, and it does not diminish
as capacity increases --- a pattern better read as a property of the imagery
than of any one model. By contrast the apparently large
\textit{Podocarpus pilgeri} difficulty is an artefact of pooling: 48 of its
55 occurrences come from ResNet-50 alone, and two architectures never make
that error. The same disaggregation removes the expectation that
Cupressaceae is the hardest family; over all five backbones it has the lowest
error rate of the five, though once ResNet-50 is excluded the single-species
Araucariaceae falls to zero and takes that position.

Neither augmentation ablation supports a recommendation. Discrete
$90^\circ$ rotation changes sign across architectures and helps neither of
the two strongest models. Removing label smoothing lowers accuracy in all
five backbones, which is the only directionally consistent pattern in the
ablation table, but each individual effect is small and every ablation was
trained with a single seed, so no effect size is established.

The most important qualification concerns what the accuracies measure. A
control experiment shows that a model trained on $16\times16$ inputs retains
95.3\,\% of full-image accuracy and one trained on heavily blurred inputs
96.6\,\%, while a colour-histogram-only model reaches 37.970\,\%. Anatomical
structure is destroyed at those resolutions, so the decisive evidence is
spatial but low-frequency, and the $99$\,\%-level accuracies cannot be read
as evidence that the models identify wood by anatomy. This does not
invalidate the architecture comparison, which was conducted under identical
conditions throughout, but it bounds what the absolute numbers may be
claimed to show. A specimen-key audit adds a second qualification: regrouping
the 552 origin keys under a coarser rule collapses them to 325, and whether
the key denotes a physical specimen rather than an image or view has not yet
been verified against the collection's records.

Three lines of work follow directly. Repeating each ablation across the same
three seeds as the main runs would make those comparisons interpretable at a
cost of 30 training runs. Completing the specimen-key verification against
the collection's own records would establish, rather than assume, that the
partition separates physical specimens. And evaluating on independently
acquired images --- different operator, equipment and lighting --- is the
only way to separate the benign and adverse readings of the probe results,
and the precondition for any claim about field use.

Until those are done, the defensible summary of this work is a controlled
comparison of five architectures on one corpus, with a clearly identified
species pair that resists every architecture tested, and an explicit warning
that the headline accuracy overstates what has been demonstrated about
anatomical identification.
